% Shared, evidence-limited report body. Only verified local PDFs support claims.

\section{Pendahuluan}

Penghitungan dan klasifikasi kematangan tandan buah segar (TBS) menuntut
deteksi yang tahan terhadap oklusi pelepah, ukuran objek yang berubah menurut
jarak, dan cahaya kanopi yang tidak seragam. RGB menyediakan warna serta
tekstur, sedangkan kedalaman dapat membantu memisahkan objek pada bidang yang
berbeda. Namun, kedalaman juga dapat berlubang, berderau, atau tidak selaras
dengan RGB; karena itu manfaat fusi harus diuji, bukan diasumsikan.

Naskah ini adalah tinjauan terverifikasi. Halaman awal PDF diperiksa terhadap
judul dan penulis pada entri. Dari 202 entri awal, 182 PDF sumber cocok dan
dipakai. Tiga PDF yang awalnya terpotong diganti dengan versi lengkap; dua PDF
yang awalnya salah dipulihkan dari berkas pengganti. Dua puluh entri lain belum
memiliki PDF yang cocok sehingga tidak dipakai untuk mendukung klaim. Kondisi
ini tidak menyatakan kedua puluh karya tersebut tidak ada, melainkan bahwa
buktinya belum tersedia di korpus lokal.

\begin{table}[htbp]
\caption{Hasil audit berkas sumber.}
\label{tab:audit}
\centering
\begin{tabular}{lr}
\toprule
Status & Jumlah \\
\midrule
PDF sumber cocok dan dipakai & 182 \\
Unduhan salah yang dipisahkan & 22 \\
Entri belum terpulihkan & 20 \\
Unduhan terpotong yang diganti versi lengkap & 3 \\
PDF tambahan tidak terhubung & 1 \\
Duplikat pengganti & 7 \\
\bottomrule
\end{tabular}
\end{table}

Tujuan laporan adalah mengidentifikasi bukti yang dapat ditransfer ke TBS tanpa
menyamakan hasil lintas domain sebagai hasil kebun. Gambar~\ref{fig:taxonomy}
menunjukkan empat komponen: deteksi RGB, kedalaman, fusi, dan aplikasi.

\begin{figure}[htbp]
\centering
\figplace{F01-taksonomi-gpt-image-2}{Taksonomi deteksi RGB, kedalaman, fusi RGB--D, dan aplikasi.}
\caption{Ruang lingkup tinjauan terverifikasi.}
\label{fig:taxonomy}
\end{figure}

\section{Metode audit dan profil korpus}

Audit memeriksa identitas dokumen, keterbacaan, dan kelengkapan halaman. Berkas
salah, tidak lengkap, tidak terhubung, serta duplikat dipisahkan dari korpus
kerja. Tabel~\ref{tab:themes} adalah distribusi 182 PDF yang identitasnya
cocok---bukan distribusi 202 entri awal. Tidak ada studi TBS RGB--D dalam
korpus yang cukup untuk dijadikan bukti performa langsung.

\begin{table}[htbp]
\caption{Distribusi 182 sumber terverifikasi menurut tema.}
\label{tab:themes}
\centering\small
\begin{tabular}{lr@{\quad}lr}
\toprule
Tema & n & Tema & n \\
\midrule
Fondasi RGB & 39 & Estimasi kedalaman & 21 \\
RGB--D SOD & 19 & Deteksi 3D & 17 \\
Segmentasi RGB--D & 15 & Pose 6D & 10 \\
Grasp robotik & 9 & Fusi multimodal & 7 \\
RGB--D SLAM & 7 & Survei YOLO & 7 \\
Dataset & 6 & Pedestrian RGB--T & 6 \\
YOLO plus RGB--D & 5 & Pertanian & 4 \\
Remote sensing & 4 & Medis & 3 \\
Industri & 3 & Total & 182 \\
\bottomrule
\end{tabular}
\end{table}

\section{Temuan deteksi RGB dan YOLO}

YOLOv1 memformulasikan deteksi sebagai satu pemetaan citra-ke-kotak dan kelas.
Pada PASCAL VOC 2007, paper asalnya melaporkan 63,4\% mAP pada 45 FPS untuk
YOLO penuh dan 155 FPS untuk Fast YOLO \cite{redmon2016yolo}. Angka ini
menjelaskan relevansinya untuk video, bukan kinerja untuk TBS. YOLOv2 dan
YOLOv3 memperluas prediksi multi-skala serta mekanisme kotak
\cite{redmon2017yolo9000,redmon2018yolov3}; YOLOv4, YOLOX, YOLOv7, dan YOLOv10
menunjukkan pentingnya desain pelatihan, penugasan positif, dan pascaproses
\cite{bochkovskiy2020yolov4,ge2021yolox,wang2023yolov7,wang2024yolov10}.

Bagi TBS, implikasinya adalah: fitur multi-skala perlu dipertahankan; latensi
harus diukur pada perangkat nyata; dan penghitungan perlu asosiasi antarb ingkai,
bukan hanya deteksi per gambar. Faster R-CNN, EfficientDet, DETR, dan Deformable
DETR menyediakan pembanding untuk akurasi, objek kecil, dan latensi
\cite{ren2017fasterrcnn,tan2020efficientdet,carion2020detr,zhu2021deformabledetr}.

\begin{figure}[htbp]
\centering
\figplace{F03-silsilah-rgb-gpt-image-2}{Silsilah detektor RGB.}
\caption{Keluarga pembanding untuk baseline TBS berbasis RGB.}
\label{fig:rgb}
\end{figure}

\section{Temuan kedalaman dan fusi RGB--D}

Kedalaman dapat diperoleh dari sensor terdaftar, stereo, atau estimasi
monokular. DPT, Depth Anything, Depth Anything V2, ZoeDepth, Metric3D, dan
Marigold memperluas pilihan saat sensor kedalaman tidak tersedia
\cite{ranftl2021dpt,yang2024depthanything,yang2024depthanythingv2,bhat2023zoedepth,yin2023metric3d,ke2024marigold}.
Peta relatif tidak boleh ditafsirkan langsung sebagai jarak metrik; pada TBS,
ia lebih aman sebagai isyarat pemisah bidang atau masukan dengan skor keandalan.

Eksperimen Ophoff dkk. menguji 28 posisi fusi dalam detektor dua cabang RGB dan
kedalaman. Pada data mereka, fusi menengah hingga akhir mengungguli fusi awal dan
RGB tunggal; besaran kenaikan AP bergantung pada dataset serta posisi fusi
\cite{ophoff2019multimodal}. Ini mendukung hipotesis awal, bukan konfigurasi
final sawit: cabang RGB dan kedalaman diproses terpisah lalu dipertukarkan pada
fitur multi-skala. D3Net menunjukkan bahwa kedalaman buruk perlu dapat ditolak
\cite{fan2020d3net}; CMX dan DFormer mendukung koreksi lintas-modal, bukan
penumpukan kanal yang buta \cite{zhang2023cmx,yin2024dformer}.

\begin{figure}[htbp]
\centering
\figplace{F04-strategi-fusi-gpt-image-2}{Strategi fusi awal, menengah, dan akhir.}
\caption{Titik fusi adalah variabel eksperimen.}
\label{fig:fusion}
\end{figure}

\section{Dari deteksi ke penghitungan TBS}

Penghitungan video membutuhkan identitas objek lintas bingkai. ORB-SLAM2,
ORB-SLAM3, dan DROID-SLAM menyediakan unsur pose kamera dan peta spasial
\cite{murartal2017orbslam2,campos2021orbslam3,teed2021droidslam}. Studi RGB--D
robotik mendukung pola hemat: lakukan deteksi 2D dahulu, kemudian gunakan
kedalaman serta segmentasi pada kandidat terpilih
\cite{elamraoui2024fusionvision,xu2024onboard}.

Studi pertanian terdekat menggunakan segmentasi instans dan fotogrametri
struktur-dari-gerak untuk lokasi buah apel; paper tersebut melaporkan F1 dari
0,816 pada deteksi 2D menjadi 0,881 pada deteksi serta lokasi 3D untuk data yang
diuji \cite{genemola2020fruit3d}. Sistem panen selada lapangan menegaskan bahwa
deteksi, geometri, dan tindakan robotik adalah tahap yang berbeda
\cite{birrell2020lettuce}. Kedua hasil adalah pelajaran desain, bukan angka
performa untuk sawit.

\begin{figure}[htbp]
\centering
\figplace{F05-pola-yolorgbd-gpt-image-2}{Pola integrasi YOLO dan RGB--D.}
\caption{Dua pola yang dapat dibandingkan pada data TBS.}
\label{fig:yolorgbd}
\end{figure}

\section{Rancangan eksperimen yang dapat diuji}

\begin{enumerate}
\item Rekam pasangan RGB--D terkalibrasi dan anotasi instans, kematangan,
      jarak, kondisi cahaya, oklusi, serta nilai kedalaman hilang.
\item Bandingkan baseline RGB-YOLO dengan fusi awal, fusi fitur menengah/akhir,
      dan mode RGB-\textit{fallback} saat kedalaman tidak andal.
\item Ukur AP per kelas dan ukuran objek, \textit{recall} oklusi, galat hitung
      absolut per pohon/lintasan, galat jarak, dan latensi ujung-ke-ujung.
\item Gunakan asosiasi berbasis tampilan, kedalaman, dan pose sebelum menjumlahkan
      tandan, lalu lakukan ablation pada titik fusi dan gerbang keandalan.
\end{enumerate}

Dataset NYU Depth V2, SUN RGB-D, dan ScanNet dapat mendukung pralatih komponen,
namun bukan pengganti data kebun \cite{silberman2012nyu,song2015sunrgbd,dai2017scannet}.

\begin{figure}[htbp]
\centering
\figplace{F07-funnel-sawit-gpt-image-2}{Corong bukti menuju validasi TBS.}
\caption{Bukti lintas domain harus diakhiri oleh pengujian kebun.}
\label{fig:funnel}
\end{figure}

\begin{figure}[htbp]
\centering
\figplace{F08-pipeline-sawit-gpt-image-2}{Pipeline TBS yang diusulkan.}
\caption{Pipeline evaluasi, bukan hasil eksperimen lapangan.}
\label{fig:pipeline}
\end{figure}

\section{Kesimpulan}

Laporan ini memakai 182 PDF yang cocok, bukan 202 entri awal. Bukti mendukung
YOLO sebagai baseline latensi rendah, fusi fitur menengah/akhir sebagai hipotesis
awal, dan pengurangan bobot kedalaman saat kualitasnya buruk. Bukti juga menunjukkan
bahwa penghitungan memerlukan asosiasi temporal serta geometri. Konfigurasi terbaik
untuk TBS RGB--D belum terbukti; dataset kebun berpasangan dan ablation terkontrol
adalah langkah berikutnya yang diperlukan.
